\scanpage{28}

pour t\deleted{\unsure{out}} hom d'algèbres $\mathcal A\longrightarrow\underline O$, le noyau $J\subset\mathcal A$ est un sous-fibré vectoriel \inserted{de rang $1$ \deleted{\unsure{(suppl. de $\underline O$ $\cdot1_{\mathcal A}$)}}} sur lequel $N_{\mathcal A/O}$ est nul, \inserted{(et inversement)} et $D=J\cdot E$ est alors un sous-fibré en droites de $E$ tel que $Q\mid D=0$. Inversement tt sous-fibré en droites de $E$ ayant cette propriété provient d'une section unique de $X'$ sur $X$. (Pour le vérifier, OPS $E=\mathcal A$, et $Q$ la norme, et c'est alors immédiat.)

4. La notion de structure de similitude sur $E$, et la distinction de celles qui sont fidèles, str. fidèles, \inserted{voire} non dégénérées, ``ne \unsure{dépend|dépendent} que de $E$ modulo tensorisation par un module inv'', donc ne \unsure{dépend|dépendent} que de l'algèbre d'Azumaya $\mathcal M$ associée,

\[
\mathcal M=\underline{\operatorname{End}}(E),
\]

de rang $4$ sur $E$. Si on part de une telle $\mathcal M$, la donnée d'une structure de similitude \deleted{relative à un Module \unsure{+}} équivaut à celle d'un sous-Module loc. \unsure{mon.|monog.|monogène} \inserted{$\Delta$} \unsure{de|du}

\[
\mathcal M/\underline O\,1_{\mathcal M}\overset{\mathrm{déf\unsure{n}}}{=}\Sigma(\mathcal M),
\]

(qui joue le rôle du Module des formes quadratiques)
% archive sheet number: 27
